% Common formatting commands.
% The book is printed in one ink, so hierarchy is carried by size, weight and
% small caps. Colour used to mark rubrics and prayer titles, and without it a
% black italic line is easily mistaken for prayer text, so rubrics drop a step
% in size and prayer titles take the bold face.

% heading face: bold small caps, letter-spaced
\newcommand{\headingfont}{\bfseries\scshape\addfontfeature{LetterSpace=6}}

% rubric / instruction text, a step down from the body. The reserve keeps the
% rubric with the first lines of what it introduces; without it a rubric can be
% left alone at the foot of a page.
\newcommand{\rubricsize}{\fontsize{8.5}{10.5}\selectfont}
\newcommand{\rubric}[1]{{\rubricsize\itshape #1}}
\newcommand{\rubricline}[1]{\par\vspace{0.6em}\rubric{#1}\par\vspace{0.3em}}
\newcommand{\instruction}[1]{\needspace{5\baselineskip}\rubricline{#1}}
% a rubric that introduces a prayer title: reserve for the rubric and for the
% title's own reserve, so the two never part company across a page break
\newcommand{\instructionheading}[1]{\needspace{11\baselineskip}\rubricline{#1}}
% a rubric that closes a block of prayer. Nothing follows that it must stay
% with, so it reserves its own line alone and may sit at the foot of a page.
\newcommand{\instructionclosing}[1]{\needspace{2\baselineskip}\rubricline{#1}}

% label over a prayer, kept with the text it introduces
\newcommand{\theotokion}{\needspace{5\baselineskip}{\headingfont Theotokion}}

% a division inside the canon of preparation: an ode number, or the kontakion
\newcommand{\canonlabel}[1]{\needspace{6\baselineskip}\par\vspace{0.8em}\begin{center}{\headingfont\normalsize #1}\end{center}}

% prayer title heading, centred
\newcommand{\prayer}[1]{\needspace{8\baselineskip}\par\vspace{0.3em}\begin{center}{\headingfont\normalsize #1}\end{center}}

\newcommand{\refrain}[1]{\par\vspace{0.4em}\textit{#1}\par\vspace{0.4em}}

% ornamental rule: line, star, line. The optional argument is the fraction of
% the text width taken by each line, so a flair can be matched to what it frames.
% It is set as one line of a paragraph rather than inside a center environment:
% a list injects its own breakpoint, which lets the rule be torn away from the
% text it should close.
\newcommand{\ornarule}[1][0.28]{%
  \par\noindent\hspace*{\fill}%
  \rule{#1\textwidth}{0.4pt}\raisebox{-2.58pt}{\,\footnotesize \ding{85}\,}\rule{#1\textwidth}{0.4pt}%
  \hspace*{\fill}\par
}

% A plume over a circle with a leaf either side, drawn upward from a rule. The
% bristles start inside the circle rather than at its crown, so they read as
% coming out of it, and the leaves sweep away from behind it. The leaf teeth are
% points on the leaf's own outline, so no stroke ends float free.
% #1 is the coordinate where it sits on the rule.
\newcommand{\plumeornament}[1]{%
  \begin{scope}[shift={(#1)}]
    \foreach \sx in {1, -1}{%
      \draw[line width=0.4pt, xscale=\sx, yshift=-0.9mm]
        (0.31mm,0.55mm) -- (1.00mm,0.79mm) -- (1.53mm,1.52mm) -- (2.40mm,1.43mm)
        -- (2.98mm,2.03mm) -- (3.89mm,1.93mm) -- (4.56mm,2.33mm) -- (5.39mm,2.41mm)
        -- (5.75mm,2.54mm) -- (5.47mm,2.28mm) -- (4.94mm,1.61mm) -- (4.21mm,1.29mm)
        -- (3.67mm,0.51mm) -- (2.80mm,0.48mm) -- (2.21mm,-0.21mm) -- (1.29mm,-0.01mm)
        -- (0.58mm,-0.27mm) -- cycle;}
    \draw[line width=0.4pt] (0,0) -- (0,1.2mm);
    \draw[line width=0.4pt] (0,2.6mm) circle (1.4mm);
    \begin{scope}[shift={(0,2.6mm)}]
      \foreach \a in {30, 45, ..., 150}{%
        \draw[line width=0.3pt] (\a:1.0mm)
          -- (\a:{1.0mm+2.3mm*(0.86+0.14*cos(\a-90))});}
    \end{scope}
  \end{scope}
}
% the same rule, printed only if the page has room for it. Used to close the
% book: a decoration must not push the run onto another page, nor be left alone
% on a page of its own, so where there is no room it is simply not printed.
\newcommand{\ornaruleiffits}[1][0.28]{%
  \ifdim\pagegoal<\maxdimen
    \ifdim\dimexpr\pagegoal-\pagetotal\relax>2.5\baselineskip
      \vspace{1em}%
      \ornarule[#1]%
    \fi
  \fi
}
